\chapter[1986 --- Cohen et al.]{1986: Stanley Cohen, Rita Levi-Montalcini}
\label{ch:1986_Cohen}

\section*{Main Contribution}

\textit{Discovered growth factors---proteins that stimulate cell growth and differentiation---opening new understanding of development, cancer, and wound healing.}

\section{Official Citation}

for their discoveries of growth factors

The 1986 Nobel Prize in Physiology or Medicine was awarded jointly to Stanley Cohen and Rita Levi-Montalcini for their discoveries of growth factors---the signaling molecules that control the growth, differentiation, and survival of cells. Their work revealed a fundamental principle of biology: that cells communicate with each other by secreting specific proteins that bind to receptors on the surface of target cells, triggering intracellular signaling cascades that regulate cell growth and behavior. This discovery opened up entirely new fields of biology and medicine, and it has had a significant impact on our understanding of development, cancer, and regenerative medicine.

\section{Background and Historical Context}

The phenomenon of cell growth and differentiation is one of the central questions of biology. During embryonic development, a single fertilized egg gives rise to billions of cells of many different types---neurons, muscle cells, blood cells, epithelial cells, and many others. This notable process of proliferation and differentiation must be precisely regulated: too much cell growth can lead to cancer, while too little can lead to developmental abnormalities. How does the embryo control the growth and differentiation of its cells?

By the mid-twentieth century, it was known that cells in the body produce and respond to various chemical signals, but the nature of these signals---and the mechanisms by which they regulate cell growth---were poorly understood. The study of nerve growth, in particular, had fascinated biologists for decades. It was known that neurons in the developing embryo extend long axons to reach their target tissues, and that this process of axon growth depends on the availability of target-derived signals. However, the identity of these signals and the mechanisms by which they promote nerve growth were unknown.

Rita Levi-Montalcini was born in Turin, Italy, in 1909, and trained in medicine at the University of Turin. She began studying the development of the nervous system in the late 1940s, at a time when the study of neuroembryology was in its infancy. Stanley Cohen was born in Brooklyn, New York, in 1922, and trained in biochemistry at Oberlin College and Washington University in St. Louis, where he joined the laboratory of Viktor Hamburger---a developmental biologist who had been studying nerve growth and differentiation for decades.

\section{The Discovery/Contribution}

\subsection{Rita Levi-Montalcini and the Discovery of Nerve Growth Factor}

Levi-Montalcini's journey to the discovery of nerve growth factor (NGF) began in the late 1940s, when she was studying the development of the peripheral nervous system in chick embryos. Working in Turin under difficult wartime conditions (she was forced to conduct research in a makeshift laboratory in her bedroom because of the disruption caused by World War II), Levi-Montalcini performed meticulous experiments in which she removed the limb buds of chick embryos and observed the effects on the developing nerves. She showed that when the target tissue (the limb bud) was removed, the neurons that would normally have innervated it failed to develop properly and died. Conversely, when she implanted an extra limb bud, the number of neurons that developed was increased.

These observations led Levi-Montalcini to hypothesize that the target tissue produces a diffusible substance that promotes the survival and growth of neurons. She called this substance ``nerve growth factor,'' and she spent the next decade searching for it. In 1954, Levi-Montalcini was invited to join Viktor Hamburger's laboratory at Washington University in St. Louis, where she began collaborating with Stanley Cohen.

Cohen, a biochemist with expertise in protein purification, was tasked with isolating and characterizing the nerve growth factor that Levi-Montalcini had postulated. Using a bioassay that measured the effect of tissue extracts on nerve fiber growth in chick embryos, Cohen and Levi-Montalcini showed that a substance present in mouse salivary glands (and in snake venom) was a potent stimulator of nerve fiber growth. Cohen purified this substance and showed that it was a protein---the nerve growth factor (NGF).

The discovery of NGF was a landmark in biology. It was the first growth factor to be identified, and it demonstrated that cells communicate with each other by secreting specific proteins that regulate their growth and behavior. Levi-Montalcini's bioassay---in which the growth of nerve fibers in chick embryos was used to detect and quantify NGF---was an elegant and sensitive tool that made the purification and characterization of NGF possible.

Levi-Montalcini and Cohen published their seminal paper on NGF in 1960, and the discovery was immediately recognized as a major advance. NGF was shown to promote the survival, growth, and differentiation of sympathetic and sensory neurons, and it was found to play a critical role in the development of the peripheral nervous system. The discovery of NGF also opened up a new field of research: the study of neurotrophic factors---the family of proteins that support the survival and growth of neurons.

\subsection{Stanley Cohen and the Discovery of Epidermal Growth Factor}

Cohen's second major discovery was the identification of epidermal growth factor (EGF), a protein that stimulates the growth and differentiation of epithelial cells. The discovery of EGF was serendipitous: while studying the effects of mouse salivary gland extracts on nerve growth, Cohen noticed that the extracts also caused premature opening of the eyelids and eruption of the incisor teeth in newborn mice. This observation led him to suspect that the salivary gland extracts contained a second growth factor---distinct from NGF---that stimulated the growth of epithelial tissues.

Cohen purified this second factor and named it epidermal growth factor (EGF). He showed that EGF was a small protein (53 amino acids) that promoted the proliferation of epithelial cells and the formation of new tissue. Cohen also discovered that EGF exerted its effects by binding to a specific receptor on the surface of target cells---the EGF receptor (EGFR). This was one of the first receptor tyrosine kinases to be identified, and its discovery provided the foundation for the understanding of signal transduction through receptor tyrosine kinases.

The discovery of EGF and the EGF receptor had immediate and significant implications for cancer biology. The EGF receptor was found to be overexpressed in many types of cancer, and mutations in the EGF receptor were shown to cause uncontrolled cell growth---a hallmark of cancer. This discovery led to the development of targeted cancer therapies, including gefitinib (Iressa) and erlotinib (Tarceva), which inhibit the EGF receptor and are used to treat non-small cell lung cancer and other malignancies.

\section{Impact on Modern Medicine}

The discoveries of Cohen and Levi-Montalcini have had a transformative impact on modern medicine. The concept of growth factors---and the signaling pathways through which they act---has become central to our understanding of development, cancer, wound healing, and regenerative medicine.

In cancer biology, the growth factor signaling pathway has proved to be of paramount importance. Many cancers are driven by the overactivation of growth factor signaling pathways---either through overproduction of growth factors, overexpression of growth factor receptors, or mutations in the signaling molecules downstream of the receptor. The recognition that cancer is often caused by the dysregulation of growth factor signaling has led to the development of a wide range of targeted cancer therapies, including monoclonal antibodies against growth factor receptors (such as trastuzumab/Herceptin for HER2-positive breast cancer) and small-molecule inhibitors of growth factor receptor kinases (such as imatinib/Gleevec for chronic myeloid leukemia).

In wound healing and tissue repair, growth factors play a critical role. Platelet-derived growth factor (PDGF), fibroblast growth factor (FGF), and transforming growth factor beta (TGF-$\beta$) are all involved in the complex process of wound healing, and recombinant growth factors are now used clinically to promote the healing of chronic wounds, burns, and surgical incisions. Becaplermin (Regranex), a recombinant form of PDGF, is approved for the treatment of diabetic foot ulcers.

In regenerative medicine, growth factors are used to promote the differentiation of stem cells into specific cell types, to enhance tissue repair and regeneration, and to develop new therapies for degenerative diseases. The use of growth factors to stimulate the regeneration of nerve tissue, bone, cartilage, and blood vessels is an active area of research, with the potential to transform the treatment of spinal cord injury, bone defects, and other conditions.

The EGF receptor pathway has been particularly important for the development of targeted cancer therapies. The recognition that many cancers are driven by overactive EGF receptor signaling has led to the development of monoclonal antibodies (such as cetuximab/Erbitux) and small-molecule inhibitors (such as gefitinib/Iressa and erlotinib/Tarceva) that block EGF receptor activity. These drugs have significantly improved survival for patients with certain types of lung cancer, colorectal cancer, and head and neck cancer.

Beyond cancer and wound healing, growth factors have been implicated in a wide range of diseases, including atherosclerosis, diabetes, Alzheimer's disease, and age-related macular degeneration. The study of growth factor signaling pathways has also provided insights into the mechanisms of aging and has suggested new strategies for promoting healthy aging and extending lifespan.

\section{Legacy}

Rita Levi-Montalcini continued her work on growth factors and neurobiology at the National Research Council in Rome and at the Institute of Neurobiology. She became an international advocate for women in science and for education in developing countries. She died in 2012 at the age of 103, but her work continues to inspire scientists around the world.

Stanley Cohen continued his work at Vanderbilt University, where he made further contributions to the study of growth factor signaling and the mechanisms of cell proliferation. He died in 2019, but his discoveries of NGF and EGF remain among a major in the history of biology.

The work of Cohen and Levi-Montalcini demonstrated that the growth and differentiation of cells are controlled by specific, identifiable molecules---the growth factors---and that these molecules act through specific receptors on the cell surface. This principle has become the foundation of modern cell biology, developmental biology, and cancer research, and it has led to the development of life-saving therapies for cancer, wound healing, and a wide range of other diseases. The discovery of growth factors is a testament to the power of curiosity-driven research and to the importance of following unexpected observations to their logical conclusions.


\section{Modern Developments}

Cohen and Levi-Montalcini's growth factor work has led to modern targeted cancer therapy. Understanding of growth factor signaling has led to drugs targeting EGFR (cetuximab), HER2 (trastuzumab), and VEGF (bevacizumab). Nerve growth factor (NGF) inhibitors (tanezumab) are being developed for chronic pain. Understanding of growth factor signaling has informed development of tyrosine kinase inhibitors (imatinib, erlotinib) that have transformed cancer treatment.


\section{Bibliography}

\begin{thebibliography}{9}
\bibitem{nobel-1986}
Nobel Prize Outreach, ``The Nobel Prize in Physiology or Medicine 1986,''\emph{NobelPrize.org}.\\
\url{https://www.nobelprize.org/prizes/medicine/1986/summary/}

\bibitem{lecture-1986}
Stanley Cohen, ``Nobel Lecture,''\emph{NobelPrize.org}. Winner-authored primary source; downloadable from the official Nobel Prize archive.\\
\url{https://www.nobelprize.org/prizes/medicine/1986/cohen/lecture/}
\end{thebibliography}
